\documentclass[11pt,a4paper]{article}

% ── Codificación y fuentes ────────────────────────────────────────────────────
\usepackage[utf8]{inputenc}
\usepackage[T1]{fontenc}
\usepackage{lmodern}
\usepackage[spanish]{babel}

% ── Geometría y espaciado ─────────────────────────────────────────────────────
\usepackage[top=2.5cm, bottom=2.5cm, left=2.8cm, right=2.8cm]{geometry}
\usepackage{setspace}
\setstretch{1.15}
\usepackage{parskip}

% ── Tablas ────────────────────────────────────────────────────────────────────
\usepackage{booktabs}
\usepackage{tabularx}
\usepackage{array}
\usepackage{longtable}
\usepackage{enumitem}

% ── Colores y cajas ───────────────────────────────────────────────────────────
\usepackage[dvipsnames]{xcolor}
\usepackage{tcolorbox}
\tcbuselibrary{skins, breakable}

% ── Cabeceras y pies de página ────────────────────────────────────────────────
\usepackage{fancyhdr}
\usepackage{lastpage}
\pagestyle{fancy}
\fancyhf{}
\fancyhead[L]{\small\textbf{Informe de Evaluación} --- Física para Bioanalistas}
\fancyhead[R]{\small daniela rivero 31435612}
\fancyfoot[C]{\small Página \thepage\ de \pageref{LastPage}}
\renewcommand{\headrulewidth}{0.4pt}

% ── Hiperenlaces ──────────────────────────────────────────────────────────────
\usepackage[colorlinks=true, linkcolor=NavyBlue, urlcolor=NavyBlue]{hyperref}

% ── Paleta de colores personalizados ─────────────────────────────────────────
\definecolor{desempeno_excelente}{HTML}{1A6B2F}
\definecolor{desempeno_bueno}{HTML}{2E5A9C}
\definecolor{desempeno_suficiente}{HTML}{8B6914}
\definecolor{desempeno_deficiente}{HTML}{C0392B}
\definecolor{desempeno_insuficiente}{HTML}{7B241C}
\definecolor{grisFondo}{HTML}{F5F5F5}
\definecolor{azulTitulo}{HTML}{1B2A6B}
\definecolor{verdeFortaleza}{HTML}{1A6B2F}
\definecolor{naranjaMejora}{HTML}{A04000}

% ── Estilos de cajas ─────────────────────────────────────────────────────────
\tcbset{
  cajaTitulo/.style={
    enhanced, breakable,
    colback=azulTitulo!8, colframe=azulTitulo,
    fonttitle=\bfseries\large, coltitle=white,
    attach boxed title to top left={yshift=-2mm, xshift=4mm},
    boxed title style={colback=azulTitulo},
    top=6pt, bottom=6pt, left=8pt, right=8pt,
  },
  cajaFortaleza/.style={
    enhanced, breakable,
    colback=verdeFortaleza!6, colframe=verdeFortaleza!60,
    fonttitle=\bfseries, coltitle=verdeFortaleza,
    top=4pt, bottom=4pt, left=8pt, right=8pt,
  },
  cajaMejora/.style={
    enhanced, breakable,
    colback=naranjaMejora!6, colframe=naranjaMejora!60,
    fonttitle=\bfseries, coltitle=naranjaMejora,
    top=4pt, bottom=4pt, left=8pt, right=8pt,
  },
  cajaRecomendacion/.style={
    enhanced, breakable,
    colback=grisFondo, colframe=gray!50,
    fonttitle=\bfseries, coltitle=black,
    top=4pt, bottom=4pt, left=8pt, right=8pt,
  },
}

% ─────────────────────────────────────────────────────────────────────────────
\begin{document}

% ══ PORTADA ══════════════════════════════════════════════════════════════════
\begin{center}
  {\Large\bfseries\color{azulTitulo} Universidad de Oriente/Núcleo Sucre --- Física para Bioanalistas}\\[4pt]
  {\large Informe de Evaluación Preliminar}\\[12pt]
  \rule{\linewidth}{1.2pt}\\[6pt]
  {\LARGE\bfseries daniela rivero 31435612}\\[4pt]
  \rule{\linewidth}{0.6pt}
\end{center}

% ══ FICHA TÉCNICA ════════════════════════════════════════════════════════════
\vspace{4pt}
\begin{tcolorbox}[cajaTitulo, title=Datos del informe]
\begin{tabular}{@{}llll@{}}
  \textbf{Estudiante:}  & daniela rivero 31435612  &
  \textbf{Fecha:}       & 2026-06-25 \\[4pt]
  \textbf{Archivo PDF:} & \texttt{daniela\_rivero\_31435612.pdf}  &
  \textbf{Páginas:}     & 5 \\
\end{tabular}
\end{tcolorbox}

% ══ RESULTADO GLOBAL ═════════════════════════════════════════════════════════
\vspace{6pt}
\begin{center}
  \begin{tcolorbox}[
    enhanced,
    colback=white, colframe=desempeno_excelente,
    width=0.62\linewidth,
    boxrule=2pt,
    halign=center, valign=center,
    top=8pt, bottom=8pt,
  ]
    \centering
    {\large\bfseries Puntaje sugerido}\\[6pt]
    {\Huge\bfseries\color{desempeno_excelente} 9.0/10}\\[6pt]
    {\large Nivel de desempeño: \textbf{\color{desempeno_excelente} Excelente}}
  \end{tcolorbox}
\end{center}

% ══ RESUMEN GENERAL ══════════════════════════════════════════════════════════
\section*{Resumen general}
El estudiante demuestra un excelente dominio de los principios de la física de fluidos aplicados al bioanálisis. Ha resuelto la mayoría de los problemas de manera correcta, aplicando las fórmulas adecuadas y realizando los cálculos con precisión. Se identificó un error conceptual menor en una parte de una pregunta, pero el desempeño general es muy sólido.

% ══ FORTALEZAS Y ÁREAS DE MEJORA ════════════════════════════════════════════
\vspace{4pt}
\begin{minipage}[t]{0.48\linewidth}
  \begin{tcolorbox}[cajaFortaleza, title={Fortalezas}]
\begin{itemize}[leftmargin=1.5em, itemsep=2pt]
  \item Fuerte comprensión de los principios fundamentales de la dinámica de fluidos (presión hidrostática, ecuación de continuidad, ecuación de Bernoulli, Ley de Poiseuille).
  \item Habilidad para aplicar correctamente las fórmulas y realizar cálculos precisos.
  \item Manejo adecuado de las unidades y conversiones.
  \item Presentación clara y organizada de las soluciones.
\end{itemize}
  \end{tcolorbox}
\end{minipage}
\hfill
\begin{minipage}[t]{0.48\linewidth}
  \begin{tcolorbox}[cajaMejora, title={Areas de mejora}]
\begin{itemize}[leftmargin=1.5em, itemsep=2pt]
  \item Revisión cuidadosa de la aplicación conceptual en problemas de varios pasos, especialmente al derivar cantidades a partir de relaciones o cocientes.
\end{itemize}
  \end{tcolorbox}
\end{minipage}

% ══ OBSERVACIONES POR PREGUNTA ═══════════════════════════════════════════════
\section*{Observaciones por pregunta}

\begin{longtable}{>{\bfseries}p{2.8cm} p{1.6cm} p{7.0cm} p{2.8cm}}
  \toprule
  \textbf{Pregunta} & \textbf{Puntaje} & \textbf{Evaluación} & \textbf{Errores detectados} \\
  \midrule
  \endhead
  \bottomrule
  \endfoot
    \textbf{Pregunta 1: Presión Hidrostática y Venosa} & 2.0/2.0 & El estudiante aplicó correctamente la fórmula de la presión hidrostática para determinar la altura mínima requerida. Los datos fueron identificados adecuadamente y el cálculo es correcto, con una aproximación razonable en el resultado final. & \textit{---} \\
    \midrule
    \textbf{Pregunta 2: Principio de Arquímedes} & 1.0/2.0 & El estudiante calculó correctamente la fuerza de empuje y el volumen de la muestra. Sin embargo, al determinar la densidad real del tejido, utilizó el peso aparente en lugar del empuje en la relación para la densidad, lo que llevó a un resultado final incorrecto para la densidad del tejido. El procedimiento para el empuje y el volumen es impecable, pero el error conceptual en la densidad final es significativo. & \textit{Error conceptual al calcular la densidad real del tejido: se utilizó el peso aparente en lugar del empuje en la relación de densidades.} \\
    \midrule
    \textbf{Pregunta 3: Hemodinámica (Continuidad y Bernoulli)} & 2.0/2.0 & El estudiante aplicó correctamente la ecuación de continuidad para calcular la velocidad en la zona estrecha y la ecuación de Bernoulli para determinar la presión en esa misma zona. Todos los cálculos y sustituciones son precisos y las unidades son consistentes. & \textit{---} \\
    \midrule
    \textbf{Pregunta 4: Ley de Poiseuille (Análisis Proporcional)} & 2.0/2.0 & El estudiante demostró una excelente comprensión del análisis proporcional de la Ley de Poiseuille. Identificó correctamente la relación de la tasa de flujo con la cuarta potencia del radio y calculó el porcentaje de caudal remanente de forma precisa. & \textit{---} \\
    \midrule
    \textbf{Pregunta 5: Flujo Viscoso y Resistencia} & 2.0/2.0 & El estudiante aplicó correctamente la Ley de Poiseuille para despejar y calcular la diferencia de presión requerida. Todas las conversiones de unidades, sustituciones y cálculos intermedios son correctos, llevando al resultado final preciso. & \textit{---} \\
    \midrule
\end{longtable}

% ══ ERRORES DETECTADOS ═══════════════════════════════════════════════════════
\section*{Análisis de errores}

\subsection*{Errores conceptuales}
\begin{itemize}[leftmargin=1.5em, itemsep=2pt]
  \item En la Pregunta 2, al calcular la densidad real del tejido, el estudiante utilizó incorrectamente el peso aparente en el denominador de la relación de densidades ($\rho$\_objeto = $\rho$\_fluido * Peso\_real / Empuje), en lugar del empuje.
\end{itemize}

\subsection*{Errores procedimentales}
\begin{itemize}[leftmargin=1.5em, itemsep=2pt]
  \item El error en el cálculo de la densidad en la Pregunta 2 es una consecuencia directa del error conceptual mencionado, no un error de cálculo aritmético per se.
\end{itemize}

% ══ RECOMENDACIONES AL ESTUDIANTE ════════════════════════════════════════════
\section*{Retroalimentación para el estudiante}
\begin{tcolorbox}[cajaRecomendacion, title={Dirigido a: daniela rivero 31435612}]
Tu desempeño en este examen es muy bueno. Demuestras un sólido entendimiento de los principios de la física de fluidos. Para seguir mejorando, te recomiendo revisar con especial atención la aplicación conceptual de las fórmulas en problemas de varios pasos, particularmente cuando se derivan propiedades a partir de relaciones. Asegúrate de que cada término en tus ecuaciones represente la magnitud física correcta (por ejemplo, la diferencia entre 'peso aparente' y 'empuje'). Continúa practicando la resolución de problemas y la organización de tus soluciones, ya que esto contribuye a la claridad y precisión de tus respuestas.
\end{tcolorbox}

% ══ NOTA PARA EL PROFESOR ════════════════════════════════════════════════════
\section*{Nota para el profesor}
\begin{tcolorbox}[
  enhanced, breakable,
  colback=yellow!8, colframe=yellow!60!black,
  fonttitle=\bfseries, coltitle=black,
  title={Observaciones para revisión manual},
  top=4pt, bottom=4pt, left=8pt, right=8pt,
]
El estudiante ha demostrado un nivel de comprensión muy alto en la mayoría de los temas evaluados. El único error significativo se encuentra en la Pregunta 2, donde, a pesar de calcular correctamente el empuje y el volumen, comete un error conceptual al usar el peso aparente en lugar del empuje para determinar la densidad real del tejido. Este error puntual no opaca el excelente desempeño general del estudiante en el resto del examen.
\end{tcolorbox}

% ══ PIE DE INFORME ════════════════════════════════════════════════════════════
\vspace{12pt}
\begin{center}
  \footnotesize\color{gray}
  Informe generado automáticamente por el sistema de evaluación con IA (gemini-2.5-flash) y soporte LaTex. \\
  Este documento es una evaluación \textbf{preliminar} para ser revisado por el profesor antes de ser entregado al 
  estudiante. \\
  Generado el 2026-06-25 \textbullet\ BIOANALISIS Sistema Académico de Evaluación.
  \end{center}

\end{document}
